% 04_first_run.tex - copy the Liard example, shorten it to one year, run it
\section{Run the Liard River example}
The Liard River example is a Raven model of the Liard River (65 subbasins in northern British Columbia and Yukon). A
hypothetical aquifer lies under three of its subbasins. Raven builds the MODFLOW~6 model of this aquifer itself, from four
inputs: the HRU outlines (polygons), the river lines, a raster of the land surface (a \emph{DEM}, digital elevation model)
and a raster of the bedrock surface.

\step{Copy the example into a work folder}
Work on a copy, so that \cmd{model/start} stays as it came:
\begin{shell}
$ cp -r $WS/model/start $WS/work
$ cd $WS/work
\end{shell}
\see no output. \cmd{ls} now shows the model files: \cmd{Liard.rvi}, \cmd{Liard.rvh}, \cmd{Liard.rvp}, \cmd{Liard.rvt},
\cmd{Liard.rvc}, \cmd{Liard.rvg} and the data folders.

\step{Shorten the run to one year}
The example simulates \V{t1:fullDays} days (\V{t1:fullYears} years). That takes about \V{t1:fullmin}~minutes. For a first
run, one year is enough. The \cmd{:Duration} line of \cmd{Liard.rvi} holds the number of days. Change 7305 to 365, with a
text editor (for example \cmd{nano Liard.rvi}) or with this command:
\term{t1_sed.txt}
Check the result:
\term{t1_duration.txt}
\see the line above. The number after \cmd{:Duration} is now 365.
\begin{trouble}
If \cmd{grep} still shows 7305, the \cmd{sed} pattern did not match. Open \cmd{Liard.rvi} in a text editor and change the
number by hand.
\end{trouble}

\step{Run Raven}
Raven takes the model name and an output folder. \cmd{Liard} is the model name: Raven reads \cmd{Liard.rvi}, \cmd{.rvh},
\cmd{.rvp}, \cmd{.rvt}, \cmd{.rvc} and, for groundwater, \cmd{Liard.rvg}. \cmd{-o out/} is the output folder.
\term{t1_run.txt}
\see Raven's progress (shortened here). The run takes about \V{t1:runsec}~seconds.
\begin{tocheck}
\begin{itemize}
\item At the end, between two lines of \cmd{=} signs, Raven prints \cmd{Exiting Gracefully: Successful
Simulation}.
\item Raven printed \cmd{Initializing groundwater model (MODFLOW 6)}.
\item Warnings are listed in \cmd{out/Raven\_errors.txt}, one per line: \V{t1:nwarn} lines in this run, as
  \cmd{grep -c WARNING out/Raven\_errors.txt} counts them. \V{t1:nwarnpoly} of them say that
  an HRU polygon does not match the HRU area of \cmd{Liard.rvh}. This is expected: the polygons are stand-ins, and
  \cmd{README.md} of the example explains why. Raven still passes all the recharge and seepage between the two models;
  only how it is spread over the cells is distorted. The other lines come from the original Liard model
  (\V{t1:nwarnchan} about channel cross-sections) and from gauges without observations in this one-year period
  (\V{t1:nwarndiag} lines naming a fit statistic that cannot be calculated without observations, and \V{t1:nwarnrep}
  lines saying that the warning before them was repeated, which Raven writes instead of the repeats). They are expected
  too.
\end{itemize}
\end{tocheck}
\begin{trouble}
\begin{itemize}
\item \cmd{Liard: command not found}: the variable \cmd{RAVEN} is empty in this terminal. Type the
  \cmd{export} lines of Step~\ref{st:unpack} again.
\item \cmd{Cannot find or read .rvh file} (or another file): you are not in \cmd{\$WS/work}. Type \cmd{cd \$WS/work}.
\item Any other stop: read the reason with \cmd{cat out/Raven\_errors.txt}.
\end{itemize}
\end{trouble}
